% ---
% titre: 9P - Figures géométriques planes
% theme: ES
% chapitre: Figures géométriques planes
% sous_chapitre:
% niveau: [9e]
% type: PLANIF
% tags: [c-figures-planes,c-angles,c-triangles,c-quadrilateres,c-constructions]
% source: latex
% ---

\header{Figures géométriques planes - Objectifs}

\def\arraystretch{1.2}
\begin{tabularx}{\linewidth}{|>{\centering\arraybackslash}p{3cm}|>{\centering\arraybackslash}p{3cm}|X|}
\hline
\cellcolor{themeColor!30} Livre & \cellcolor{themeColor!30} Fiches & \cellcolor{themeColor!30} \hfill Aide-mémoire\hfill\null\\
\hline
\multirow{3}{*}{pp. 105-127}&\multirow{3}{*}{pp. 103-144}&Figures géométriques planes : pp. 104-115\\
&&Angles et droites remarquables : pp. 115-137\\
&&Quadrilatères : pp. 138-139\\
\hline
\end{tabularx}
\def\arraystretch{1}


\def\arraystretch{1.6}
\begin{tabularx}{\textwidth}{|>{\raggedright\arraybackslash}X|>{\raggedright\arraybackslash}X|}
\hline
\cellcolor{themeColor!30}Compétences visées & \cellcolor{themeColor!30}Utilité dans la vie de tous les jours \\
\hline
Mobiliser les savoirs, savoir-faire et savoir être pour résoudre des problèmes en lien avec les figures géométriques.
&\vspace{-1em}
\begin{itemize}[leftmargin=*, itemsep=2pt, topsep=0pt, partopsep=0pt, parsep=0pt, label=$\triangleright$]
    \item Construction d'un meuble
    \item Agencement d'un appartement
    \item Dessin
\end{itemize} \\
\hline

Développer des stratégies de recherche afin de résoudre des problèmes.
&\vspace{-1em}
\begin{itemize}[leftmargin=*, itemsep=2pt, topsep=0pt, label=$\triangleright$]
    \item Découper un projet en sous-tâches
    \item Choisir l'ordre des tâches à effectuer
\end{itemize} \\
\hline

Présenter sa démarche de façon claire (explications, grandeur cherchée, unités, calculs, réponse).
&\vspace{-1em}
\begin{itemize}[leftmargin=*, itemsep=2pt, topsep=0pt, label=$\triangleright$]
    \item Collaborer avec d'autres personnes
    \item Expliciter sa compréhension d'un sujet
    \item Argumenter
\end{itemize} \\
\hline

Développer des stratégies d'apprentissage pour s'approprier puis mobiliser les savoirs et savoir-faire.
&\vspace{-1em}
\begin{itemize}[leftmargin=*, itemsep=2pt, topsep=0pt, label=$\triangleright$]
    \item Apprendre une nouvelle activité %(art, sport, \ldots)
    \item S'adapter à un nouvel environnement %(travail, voyage, \ldots)
    \item Se former en autonomie %(tutoriel, lecture, \ldots)
\end{itemize} \\
\hline

Vérifier sa solution ou au moins sa plausibilité.
&\vspace{-1em}
\begin{itemize}[leftmargin=*, itemsep=2pt, topsep=0pt, label=$\triangleright$]
    \item S'autocritiquer
    \item Évaluer la qualité de son travail
\end{itemize} \\
\hline
\end{tabularx}
\def\arraystretch{1}


\vfill
\begin{astuce}{Ce que je dois déjà savoir}
\begin{itemize}[leftmargin=*, itemsep=2pt, topsep=2pt, label=$\triangleright$]
    \item Utiliser correctement la règle, l'équerre et le compas
\end{itemize}
\end{astuce}

\newpage
%---------------
\def\arraystretch{1.6}
\begin{tabularx}{\textwidth}{|>{\raggedright\arraybackslash}p{5.2cm}|>{\raggedright\arraybackslash}X|>{\centering\arraybackslash}p{0.8cm}|>{\centering\arraybackslash}p{0.8cm}|>{\centering\arraybackslash}p{0.8cm}|}
\hline
\rowcolor{themeColor!30}\multicolumn{2}{|l|}{Objectifs} &\multicolumn{3}{c|}{Auto-évaluation} \\
\rowcolor{themeColor!30}\multicolumn{2}{|l|}{} & \Sadey[1.5]& \Neutrey[1.5] & \Smiley[1.5]  \\\hline

\multirow{5}{=}{Reconnaître, nommer, décrire des figures planes selon leurs propriétés (symétries internes, côtés, angles, somme des angles, diagonales) et construire des angles, triangles, quadrilatères, cercles.}
& Connaître le vocabulaire et les concepts liés aux figures planes :

droite, demi-droite, segment, distance entre deux points, entre un point et une droite, entre deux droites parallèles, droites parallèles, sécantes, perpendiculaires, côté, sommet, base, hauteur, longueur, largeur, grande et petite base, diagonale, cercle, arc de cercle, centre, corde, rayon, diamètre, axe et centre de symétrie, périmètre, aire. & & & \\
\cline{2-5}
& Connaître la définition et les propriétés des triangles et quadrilatères particuliers pour pouvoir les identifier, les classer, les construire ou déduire certaines informations manquantes :
\begin{itemize}[leftmargin=*, itemsep=0pt, topsep=2pt, partopsep=0pt, parsep=0pt, label=--]
    \item triangles : quelconque, rectangle, isocèle, équilatéral, rectangle isocèle
    \item quadrilatères : trapèze, trapèze rectangle et isocèle, parallélogramme, carré, losange, rectangle
\end{itemize} & & & \\
\cline{2-5}
& Noter un angle ou pouvoir interpréter sa notation. & & & \\
\cline{2-5}
& Construire un triangle en connaissant ses trois côtés, deux côtés et un angle, un côté et deux angles. & & & \\
\cline{2-5}
& Construire un quadrilatère en fonction de ses propriétés et en connaissant certaines de ses mesures (côtés, angle, diagonale). & & & \\
\hline

\multirow{1}{=}{Reconnaître et nommer des angles (aigu, obtus, droit, plat, rentrant).}
& Donner le type d'un angle en fonction de sa mesure ou pouvoir le représenter en fonction de son type : angle nul, aigu, droit, obtus, plat, rentrant, plein. & & & \\
\hline
\end{tabularx}

\newpage

\begin{tabularx}{\textwidth}{|>{\raggedright\arraybackslash}p{5.2cm}|>{\raggedright\arraybackslash}X|>{\centering\arraybackslash}p{0.8cm}|>{\centering\arraybackslash}p{0.8cm}|>{\centering\arraybackslash}p{0.8cm}|}
\hline

\multirow{3}{=}{Estimer, comparer, classer et mesurer des angles en degrés.}
& Estimer à vue d'œil la valeur d'un angle en reconnaissant son type. & & & \\
\cline{2-5}
& Construire et mesurer des angles avec un rapporteur. & & & \\
\cline{2-5}
& Faire des calculs d'angles en utilisant :
\begin{itemize}[leftmargin=*, itemsep=0pt, topsep=2pt, partopsep=0pt, parsep=0pt, label=--]
    \item les définitions de l'angle aigu, obtus, rentrant, plein, plat, droit
    \item les propriétés de la somme des angles d'un triangle et d'un quadrilatère
    \item les propriétés des polygones
\end{itemize} & & & \\
\hline

\multirow{1}{=}{Reconnaître, nommer, décrire des propriétés et construire des droites parallèles et perpendiculaires.}
& Construire des droites perpendiculaires et parallèles de manières différentes :
\begin{itemize}[leftmargin=*, itemsep=0pt, topsep=2pt, partopsep=0pt, parsep=0pt, label=--]
    \item à l'aide d'une équerre et d'une règle
    \item à l'aide d'un compas et d'une règle
\end{itemize} & & & \\
\hline

\multirow{2}{=}{Représenter des figures planes par un croquis et/ou un dessin à l'échelle (y compris 1:1).}
& Différencier un croquis d'un dessin à l'échelle. & & & \\
\cline{2-5}
& Être capable de réaliser un croquis et de le coder ou d'interpréter un codage. & & & \\
\hline

\end{tabularx}
\def\arraystretch{1}

\vfill
\begin{astuce}{Mon bilan}
Ce que je dois encore travailler~:\\[2pt]
\noindent\makebox[\linewidth]{\dotfill}\\[2pt]
\noindent\makebox[\linewidth]{\dotfill}\\[7pt]
Ce que je peux faire pour progresser~:\\[2pt]
\noindent\makebox[\linewidth]{\dotfill}\\[2pt]
\noindent\makebox[\linewidth]{\dotfill}
\end{astuce}

%------------------------------------------------
\newpage
\header{Figures géométriques planes - Feuille de route}

\newcommand{\rref}[1]{\textcolor{themeColor}{#1}}

% Case à cocher
\newcommand{\case}{\raisebox{-0.1ex}{\Large$\square$}}

% Marqueurs obligatoire / facultatif (forme, pas seulement couleur)
\newcommand{\oblig}{\tikz[baseline=-0.5ex]{\fill[themeColor] (0,0) circle (0.09);}}
\newcommand{\faclt}{\tikz[baseline=-0.5ex]{\draw[gray!70,line width=0.6pt] (0,0) circle (0.09);}}

\vfill
\begin{tikzpicture}[
circ/.style={circle, fill=themeColor, text=white, font=\bfseries\Large, minimum size=11mm, inner sep=0pt},
 box/.style={rectangle, draw=themeColor!35!white, rounded corners=2pt, fill=themeColor!8!white, inner sep=3mm, text width=14.7cm, align=left, anchor=north west}]

% ---------- Étape 1 ----------
\node[circ, anchor=north] (c1) at (0,0) {1};
\node[box] (b1) at ($(c1.north)+(1.2cm,0)$) {%
\textbf{\large Propriétés des quadrilatères}\\[4pt]
{\small\begin{tabularx}{10cm}{cXp{3cm}p{1cm}}
\case & \rref{Que Sais-Je ?} & Fiche, pp.\,103-105 & \oblig \\
\case & \rref{ES 75} & Fiche, p.\,133 & \oblig \\
\case & \rref{ES 78} & Livre, p.\,124 & \oblig \\
\case & \rref{ES 82} & Livre, p.\,124 & \oblig \\
\end{tabularx}}
};

% ---------- Étape 2 ----------
\coordinate (y2) at ($(b1.south west)+(0,-3mm)$);
\node[circ, anchor=north] (c2) at (c1 |- y2) {2};
\node[box] (b2) at (y2) {%
\textbf{\large Nomenclature et mesures d'angles}\\[4pt]
{\small\begin{tabularx}{10cm}{cXp{3cm}p{1cm}}
\case & \rref{Constellations} & Fiche annexe & \oblig \\
\end{tabularx}}
};

% ---------- Étape 3 ----------
\coordinate (y3) at ($(b2.south west)+(0,-3mm)$);
\node[circ, anchor=north] (c3) at (c1 |- y3) {3};
\node[box] (b3) at (y3) {%
\textbf{\large Construction des quadrilatères}\\[4pt]
{\small\begin{tabularx}{10cm}{cXp{3cm}p{1cm}}
\case & \rref{ES 83} & Livre, p.\,124 & \oblig \\
\case & \rref{ES 84} & Livre, p.\,125 & \oblig \\
\case & \rref{ES 85} & Livre, p.\,125 & \oblig \\
\case & \rref{ES 86} & Livre, p.\,125 & \oblig \\
\case & \rref{ES 87} & Livre, p.\,125 & \oblig \\
\end{tabularx}}
};

% ---------- Étape 4 ----------
\coordinate (y4) at ($(b3.south west)+(0,-3mm)$);
\node[circ, anchor=north] (c4) at (c1 |- y4) {4};
\node[box] (b4) at (y4) {%
\textbf{\large Calculs d'angles}\\[4pt]
{\small\begin{tabularx}{10cm}{cXp{3cm}p{1cm}}
\case & \rref{Dossier Biceps} & Dossier annexe & \oblig \\
\case & \rref{ES 94} & Livre, p.\,126 & \oblig \\
\case & \rref{Problème 109} & VLM, p.\,95 & \oblig \\
\end{tabularx}}
};

% ---------- Étape 5 ----------
\coordinate (y5) at ($(b4.south west)+(0,-3mm)$);
\node[circ, anchor=north] (c5) at (c1 |- y5) {5};
\node[box] (b5) at (y5) {%
\textbf{\large Exercices \guillemotleft\,Défi\,\guillemotright}\\[4pt]
{\small\begin{tabularx}{10cm}{cXp{3cm}p{1cm}}
\case & Problème 74 & VLM, p.\,66 & \faclt \\
\case & ES 89 & Livre, p.\,126 & \faclt \\
\case & Problème 110 A. & VLM, p.\,95 & \faclt \\
\end{tabularx}}
};

% ---------- Ligne de parcours (derrière les cercles) ----------
\begin{pgfonlayer}{background}
\draw[themeColor, line width=2pt] (c1) -- (c2) -- (c3) -- (c4) -- (c5);
\end{pgfonlayer}

\end{tikzpicture}

\vfill
\begin{center}
\small \case~à faire \qquad \oblig~\rref{obligatoire} \qquad \faclt~facultatif (consolidation / pour aller plus loin)
\end{center}
\vspace{1em}

%------------------------------------------------